% ============================================================================
% Section 4: Experimental Design
% Sources: writeup/contents/methods.tex (auctions, SMAD, reconstruction, models)
%          Engineering_simplicity/main.tex:126-274 (DA, prompt skeletons, models)
% Per plan/plan.md section 3 row 4 and section 5 reconciliation decisions.
% ============================================================================

\section{Experimental Design}\label{sec:design}

Our design crosses the lever families of Section~\ref{sec:framework} with the
canonical strategy-proof auction domain---second-price sealed bidding and its
clock implementations---and four LLM families, plus one field-inspired
environment (a stylized eBay marketplace) that serves validation only. A second
strategy-proof domain, school choice under student-proposing deferred
acceptance, replicates the design end-to-end and serves as a generalization
check on the ranking; its design is summarized below
(Section~\ref{sec:design-da}) and its implementation and results are reported
in Appendix~\ref{app:da}. The unit
of observation throughout is a fully specified, seed-reproducible \emph{market
instance}: a draw of values (or preferences), a mechanism, a description, and a
set of LLM agents whose actions mechanically determine allocations and
payments. This section describes the mechanisms and environments
(Section~\ref{sec:design-environments}), the treatment grid
(Section~\ref{sec:design-treatments}), the models and simulation protocol
(Section~\ref{sec:design-models}), the unified deviation metrics and the
preservation index that aggregates them into the ranking
(Section~\ref{sec:metrics}), and the reconstruction of human benchmarks
(Section~\ref{sec:reconstruction}). Table~\ref{tab:calibration-master} pins the
canonical calibration that all cells share; where legacy runs deviate from this
calibration, the deviation is flagged in the table notes and resolved by the
harmonized re-runs.

\subsection{Mechanisms and Environments}\label{sec:design-environments}

\subsubsection{Auctions}\label{sec:design-auctions}

We generate synthetic auction data by simulating auctions in which $n$ LLM
agents bid against one another under fully specified information structures and
auction rules. Each experimental \emph{setting} is defined by (i) a \emph{value
model} (how bidder values or signals are drawn), (ii) a \emph{mechanism
environment} (the action space, allocation and payment rules, and the
information revealed during play), and (iii) a \emph{prompt framing} (how the
mechanism is described to the bidder). For each setting, we repeat the
simulation $K$ times with independent random seeds, producing a dataset of
realized values (or signals), bidder actions, allocations, prices, profits, and
free-text plans. ``Actions'' refers to: a numeric bid in sealed-bid formats, a
\texttt{stay}/\texttt{exit} decision at each clock tick in ascending auctions,
and an \texttt{increase}/\texttt{hold} decision over a bidder's maximum bid in
the eBay proxy-bidding environment.

\paragraph{Value models and numerical calibration.}
We implement three canonical information structures. All parameters (support
ranges, bidder count $n$, and bid/price increments) are specified in
configuration files and held fixed within a setting. In sealed-bid auctions,
bids are restricted to a discrete grid with increment $\$0.1$; in ascending
clock auctions, the clock price increases by $\$0.5$ per
tick.\footnote{The predecessor draft of the ranking study stated a $\$0.01$
sealed-bid increment; the configuration files specify $\$0.1$, which we adopt
as canonical here (Table~\ref{tab:calibration-master}).}

\begin{itemize}
    \item \textbf{Independent private values (IPV).} Each bidder's value is an
    independent private draw from the integer grid
    \[
          v_{i} \sim \mathrm{Unif}\{0,1,\dots,49\},
    \]
    so that the expected truthful bid in the second-price auction is
    $\mathbb{E}[v]=24.5$. This is the baseline environment for the lever grid.

    \item \textbf{Affiliated private values (APV).} Values share a common
    component plus an idiosyncratic private shock:
    \[
        c \sim \mathrm{Unif}[0,29], \qquad \epsilon_{i} \sim \mathrm{Unif}[0,20],
        \qquad v_{i}=c+\epsilon_{i}.
    \]
    This induces affiliation: higher realized values for one bidder make higher
    values for others more likely. Each bidder observes her own realized
    $v_{i}$, so the dominant strategy in strategy-proof formats is unchanged.

    \item \textbf{Common value (CV).} The item has an underlying common value
    $V$ shared by all bidders, while each bidder observes a noisy private
    signal:
    \[
        V \sim \mathrm{Unif}[20,29], \qquad \eta_{i} \sim \mathrm{Unif}[-20,20],
        \qquad s_{i}=V+\eta_{i}.
    \]
    Bidder $i$ conditions on $s_{i}$ when bidding, while realized payoffs
    depend on $V$, generating the winner's-curse selection effect.
\end{itemize}

\paragraph{Sealed-bid auctions.}
In sealed-bid auctions, each bidder submits a single bid $b_{i}$
simultaneously. Let $b_{(1)}\ge b_{(2)}\ge b_{(3)}$ denote the order statistics
and let $i^\star$ be the highest bidder. The winner pays a price $p$ determined
by the mechanism:
\begin{itemize}
    \item \textbf{First-price sealed-bid (FPSB).} $i^\star$ wins and pays
    $p=b_{(1)}$.
    \item \textbf{Second-price sealed-bid (SPSB).} $i^\star$ wins and pays
    $p=b_{(2)}$.
    \item \textbf{Third-price sealed-bid (TPSB).} $i^\star$ wins and pays
    $p=b_{(3)}$.
\end{itemize}
Profits are computed as $\pi_{i}=v_{i}-p$ for winners (or $\pi_{i}=V-p$ in CV)
and $0$ for losers. SPSB is the strategy-proof baseline on which the lever grid
operates; FPSB and TPSB appear only in the behavioral-fidelity battery of
Section~\ref{sec:validation}, where their Bayes--Nash benchmarks
(Section~\ref{sec:metrics}) discipline the comparison with human
data.\footnote{The repository also contains FPSB and TPSB \emph{intervention}
batteries; these are outside the scope of the ranking, which concerns
strategy-proof mechanisms whose benchmark is a dominant strategy rather than a
Bayes--Nash equilibrium. They enter the paper instead through the
cardinal-calibration analysis of Section~\ref{sec:humans}, where the
non-strategy-proof formats discipline the search for prompt-level tunings that
match human levels and directions.}

\paragraph{Ascending clock auctions.}
In an ascending clock auction, the price starts at $0$ and increases by a fixed
increment $\Delta_{\text{clock}}=\$0.5$ at each tick. At each posted price,
each active bidder chooses \emph{stay} or \emph{exit}; exit is irreversible.
The last remaining bidder wins and pays the exit price of the second-to-last
bidder, implementing the standard second-price logic in a dynamic interface. We
distinguish two information regimes: in the \emph{open} clock (AC), bidders are
informed as rivals drop out; in the \emph{closed} or \emph{blind} clock (AC-B),
exit decisions are private and bidders are never told when others leave. Both
regimes preserve the dominant strategy of remaining active exactly while the
price is below one's value. The clock format imported from the ranking
pipeline is AC-B.

\paragraph{Why the sealed-bid-versus-clock comparison runs in both value
environments.}
The extensive-form lever (Family~A) appears twice in the auction design---once
under APV and once under IPV---and the two cells do different work. The
\emph{APV cells exist for human-benchmark comparability}: the laboratory
evidence that clock formats improve truthful play
\citep{li2017obviously,breitmoser2022obviousness} is the evidence base for our
moment-matched human reconstructions (Section~\ref{sec:reconstruction}), and it
is in this environment that humans and LLMs can be placed on the same SMAD
scale across SP, AC, and AC-B. The \emph{IPV cells exist for identification}:
under independent private values, a bidder who knows her own value learns
nothing payoff-relevant from the clock---no information about her value exists
in others' dropout decisions---so any improvement from the clock format must
operate through a purely cognitive channel. In a common- or affiliated-value
setting, by contrast, an open clock could in principle improve play by
revealing information about others' values, confounding the cognitive
interpretation. Keeping both cells, purpose-labeled, lets the APV cell carry
the human anchor and the IPV cell carry the causal claim.
The imported-pipeline clock cells are APV (common component $[0,29]$) with open
and blind variants and mixed increments (\$1 vs.\ \$0.5)---not the IPV clock the
predecessor draft describes---so the clean IPV-clock identification cell is run
separately (IPV template, IPV draws, tick \$0.5, $K=50$). Two of the four
canonical cells are complete---GPT-4o (SMAD $1.5\%$) and Gemma (SMAD $1.8\%$),
folded into Section~\ref{sec:ranking}. \TODO{IPV clock for Claude 3.5 Haiku and
Gemini 2.0 Flash pending native keys (neither is servable on OpenRouter, E2).}

\paragraph{Stylized eBay proxy-bidding environment (validation only).}
Separately from the laboratory mechanisms, we implement a stylized eBay
environment with \emph{proxy bidding}: each bidder maintains a maximum bid (a
willingness-to-pay cap), and the platform automatically raises the standing bid
only as needed to keep the current leader ahead. Each eBay simulation has $n=3$
bidders, IPV values $v_i \sim \mathrm{Unif}[0,99]$, and $T_{\text{eBay}}=10$
discrete bidding periods; treatments cross a hard versus soft (auction
extension) closing rule with the presence of a hidden reserve price. Because
this environment studies a richer, marketplace-inspired microstructure rather
than a canonical laboratory mechanism, it enters the paper only as a field
validation exhibit (Section~\ref{sec:validation}); the eBay prompts are
deliberately minimal and descriptive, so that behavior is not steered by
didactic instruction. Full rules, treatments, and procedures appear in
Appendix~\ref{app:ebay}.

\subsubsection{Deferred acceptance (generalization domain)}\label{sec:design-da}

The matching domain replicates the auction design end-to-end in a school-choice
environment in which LLMs play the role of students; because the paper's main
domain is auctions, we summarize the design here and defer the full
implementation---priority orders, prompt skeletons, and the query
algorithm---to Appendix~\ref{app:da}. Each market instance comprises four
students and four schools with unit capacity: large enough to admit meaningful
strategic considerations (competition for popular schools and priority
asymmetries), small enough that the full DA trace can be logged and every
outcome audited. Valuations are affiliated, $v_{i,s} = c_s +
\varepsilon_{i,s}$, with a school-level common component $c_s \sim
\mathrm{Unif}[40,70]$ and independent student--school shocks
$\varepsilon_{i,s} \sim \mathrm{Unif}[0,20]$ \citep{klijn2019static}; a
student's \emph{true preference ranking} is the descending sort of her
valuations. Schools rank students by fixed, Ergin-acyclic priority orders
\citep{ergin2002efficient}---acyclicity being the condition under which a
sequential-query implementation of DA can be made obviously strategy-proof
\citep{ashlagi2018stable}. All students additionally observe a fixed ``global
popularity'' ranking that proxies common beliefs about demand. Because DA is
strategy-proof for the proposing side \citep{roth1982economics}, beliefs about
others' demand should not change any report; the fixed signal exists to give
belief-directed misreporting room to express itself---and to give the belief
scaffolds of Family~C3 something concrete to amplify.

Each instance is elicited under two protocols, mirroring the auction domain's
sealed-bid/clock pair. Under \emph{direct DA}---the strategy-proof
baseline---each student submits a complete rank-order list in a single shot,
and student-proposing DA runs on the submitted rankings with the true school
priorities. Under \emph{iterative DA}---the Family~A lever in matching---no
full ranking is ever submitted: the mechanism queries students sequentially
with yes/no offers over individual schools, following the OSP tree
construction of \citet{ashlagi2018stable}, and reduces to serial dictatorship
(a single pick from the remaining schools) when only one top-priority student
remains. Within each repetition we query one LLM instance per student with a
self-contained scenario prompt bundling the student's identity, priorities,
the popularity signal, the mechanism rules, and any treatment text, with
chain-of-thought \citep{wei2022chain} elicited through structured reasoning
and decision tags. For each treatment condition and mechanism we run $K=50$
independent repetitions (fresh value draws; priorities and popularity signal
fixed), logging a complete JSON record per instance---raw responses, parsed
decisions, and the full DA trace---for exact replay and audit. The priority
orders, both prompt skeletons, and Algorithm~1 (the query protocol) appear in
Appendix~\ref{app:da} and Appendix~\ref{app:prompts}.

% ---------------------------------------------------------------------------
% Master calibration table (plan section 5.4)
% ---------------------------------------------------------------------------
\begin{table}[t]
\centering
\caption{Master calibration. Canonical parameterization of every environment
in the design; all cells use temperature $T=0.5$ and a one-shot protocol (no
feedback across instances). $K$ = repetitions per cell (pinned targets;
deviations of legacy runs flagged in the notes).}
\label{tab:calibration-master}
\footnotesize
\setlength{\tabcolsep}{3.5pt}
\resizebox{\textwidth}{!}{%
\begin{tabular}{@{}lllllcl@{}}
\toprule
Environment & Mechanisms & Value model & Action grid & $N$ & $K$ per cell & Model coverage \\
\midrule
Sealed bid, IPV & SPSB (anchor); FPSB, TPSB & $v_i \sim \mathrm{Unif}\{0,\dots,49\}$ & bid grid \$0.1 & 3 & 100 / $\ge$50$^{a}$ & canonical four + tiers$^{b}$ \\
Sealed bid, APV & SP & $v_i = c+\epsilon_i$;\; $c\sim U[0,29]$, $\epsilon_i\sim U[0,20]$ & bid grid \$0.1 & 3 & $\ge$50 & canonical four + tiers$^{b}$ \\
Clock, APV & AC, AC-B & as APV above & tick \$0.5 & 3 & 50$^{c}$ & canonical four + tiers$^{b}$ \\
Clock, IPV & AC-B & as IPV above & tick \$0.5 & 3 & 50$^{d}$ & canonical four$^{d}$ \\
Sealed bid, CV & FPSB, SPSB & $V\sim U[20,29]$;\; $s_i = V+\eta_i$, $\eta_i\sim U[-20,20]$ & bid grid \$0.1 & 4--7 & ---$^{e}$ & GPT-4o \\
eBay (validation) & proxy bidding; closing rule $\times$ reserve & $v_i \sim U[0,99]$ & raise/hold max bid & 3 & 30$^{e}$ & GPT-4o \\
Matching (DA)$^{f}$ & direct; iterative (OSP) & $v_{i,s}=c_s+\varepsilon_{i,s}$;\; $c_s\sim U[40,70]$, $\varepsilon_{i,s}\sim U[0,20]$ & rank-order list / yes--no & $4\times4$ & 50 & canonical four \\
\bottomrule
\end{tabular}}
\par\smallskip
\begin{minipage}{\textwidth}\footnotesize
\emph{Notes.} $^{a}$\,$K=100$ for the anchor cell (SPSB-IPV baseline, per
model); $K\ge 50$ for lever-grid and robustness cells.
\TODO{E8: harmonized baseline re-run pending; some legacy robustness cells ran
$K=30$--$40$.}
$^{b}$\,Canonical four = GPT-4o (2024-08-06, anchor), Claude 3.5 Haiku
(2024-10-22), Gemini 2.0 Flash, Gemma 27B; appendix tiers = Llama-3-8B
(capability floor) and the frontier tier gpt-5-mini, gpt-5, claude-sonnet-5,
and gemini-2.5-flash (full lever grid; the three servable of these also run
the DA grid, gpt-5 excepted), Appendix~\ref{app:ablations}. Gemma ran the
six-format sealed-bid fidelity battery (Table~\ref{tab:gemma-fidelity}); its
two ascending-clock cells are pending native-key runs.
$^{c}$\,Clock $K$: the newer clock cells (the harmonized IPV-clock cells and
the frontier-battery AC-B cells) run the pinned $K=50$; some legacy incumbent
APV clock cells still stand at $K=40$ (configs \texttt{repetitions: 40}) and
will be topped up in the harmonized re-runs.
$^{d}$\,IPV clock (identification cell): the clean $K=50$ cells now exist for
GPT-4o (SMAD $1.5\%$) and Gemma (SMAD $1.8\%$, $K=44$ after parse drops) via
OpenRouter, folded into Section~\ref{sec:ranking}; Claude 3.5 Haiku and Gemini
2.0 Flash are not servable on OpenRouter and their two IPV-clock cells await
native keys. (The legacy imported clock cells are APV with mixed increments,
\$1 vs.\ \$0.5, not the IPV clock described in the predecessor draft.)
$^{e}$\,eBay: $K=30$ runs per treatment cell, verified from the recovered
logs (29 for the soft-close $r=60$ cell; Appendix~\ref{app:ebay}).
\TODO{Pin $K$ for the CV cells from the recovered run logs (E5/E10).}
$^{f}$\,Generalization domain: design summary in
Section~\ref{sec:design-da}; implementation details, Algorithm~1, and results
in Appendix~\ref{app:da}.
\end{minipage}
\end{table}

\subsection{Treatments}\label{sec:design-treatments}

The treatment grid crosses the lever families of Section~\ref{sec:framework}
with the two strategy-proof domains. Family~A (extensive form) is implemented
as a change of \emph{mechanism}: the ascending clock replaces the sealed-bid
interface in auctions, and the iterative query protocol replaces direct
revelation in DA, with value draws, allocation and payment rules, model, and
temperature all held fixed. Families~B and~C are implemented as \emph{prompt
treatments} on the sealed-bid and direct-DA baselines: each treatment modifies
only the rule-explanation block of the prompt while holding the value
distribution, bidder identities, payment rules, and the structured response
format fixed, so that any behavioral change is attributable to the description
or scaffold alone, under the same underlying game.

\paragraph{Family B: descriptions.}
Three sub-families restate or expose properties of the mechanism without
changing its extensive form.
\begin{itemize}
    \item \textbf{B1, menu restatement.} Bidders are instructed to choose a
    ``maximum acceptable price'': they win if the highest rival bid (the market
    price) is below their chosen maximum, and if they win they pay that market
    price. This is the menu-style description of the second-price auction
    studied by \citet{gonczarowski2023strategyproofness}. The DA analogue
    restates DA through the menu of schools a student's report can obtain
    given others' reports, and is implemented in two variants---a
    \emph{property} variant that includes the menu's invariance statement (the
    student's report cannot change which schools are obtainable) and a
    \emph{mechanics} variant that restates the computation only---so that the
    content of a restatement can be separated from its form
    (Appendix~\ref{app:da}, Appendix~\ref{app:prompts}). This split turns out
    to be load-bearing for the ranking's interpretation
    (Section~\ref{sec:ranking-synthesis}).
    \item \textbf{B2, safety-exposing descriptions.} One-sentence statements of
    the incentive invariant that makes truthful play safe: \emph{Payoff Safety}
    in auctions (``your bid determines \emph{if} you win, not \emph{what} you
    pay'') and \emph{Rejection Safety} in DA (``rejections only redirect, never
    eliminate''). Exact texts appear in Section~\ref{sec:ranking-descriptions}
    and Appendix~\ref{app:prompts}.
    \item \textbf{B3, clock-framing.} The sealed bid is \emph{described} as an
    automatic exit threshold in a simulated ascending clock: conceptually, the
    clock rises until bidders ``drop out'' at their thresholds, and the winner
    pays the second-highest threshold. Both B1 and B3 implement exactly the
    same allocation and payment rule as the SPSB auction; only the prompt
    language differs. B3 is a description, \emph{not} a mechanism change, and
    must never be conflated with the true clock formats AC/AC-B of
    Family~A.\footnote{\TODO{The repository template name
    \texttt{intervention\_proxy\_breitmoser} historically referred to both the
    clock-framing description and the true iterative clock; prompt IDs are
    split into \texttt{desc\_clock\_framing} (B3) and
    \texttt{osp\_clock\_iterative} (A) per plan \S5.6, and the provenance of
    the imported clock cells is being audited.}}
\end{itemize}

\paragraph{Family C: reasoning scaffolds.}
Three sub-families scaffold a mode of reasoning without revealing the dominant
strategy. \textbf{C1, contingent reasoning} (the \emph{Payoff Tree}): the
prompt lays out the explicit decision tree of outcomes---win versus lose paths
in auctions, accept versus reject branches in DA---telling the agent \emph{what
happens} under each action, not what to do; run in both domains. \textbf{C2,
forward planning} (one- and two-step \emph{lookahead}): the prompt asks the
agent to simulate the algorithm's future rounds before acting. C2 runs in DA
only: DA's underlying algorithm processes the submitted ranking dynamically
through rounds of proposals and rejections, giving forward planning a natural
place to bind, whereas the sealed-bid auction has no analogous internal
dynamics. \textbf{C3, belief formation} (first- and second-order beliefs): the
prompt directs attention to what opponents will do, or to what opponents think
the agent will do; run in both domains.

\paragraph{Family D: preference framings.}
Prospect-theory-inspired personas (risk aversion, loss aversion, endowment
framing) alter the agent's stated objective rather than the mechanism's
presentation. They are not designer levers in the same sense---no designer
would instruct participants to be risk-averse---but they bound the ranking from
below and are reported in Appendix~\ref{app:prospect}.

\paragraph{Inventory and coverage.}
Every treatment template, its exact prompt text, its mapping from legacy
template names to the taxonomy above, and its run status appear in
Appendix~\ref{app:prompts} (Table~\ref{tab:intervention-inventory}). The
menu-restatement (B1) and clock-framing (B3) auction cells, GPT-4o-only in the
predecessor draft, now cover all four canonical models at $K \approx 50$
(Section~\ref{sec:ranking-descriptions}), as do the DA menu variants
(Appendix~\ref{app:da}).

\subsection{Models and Protocol}\label{sec:design-models}

\paragraph{Models.}
The canonical model set comprises four widely deployed, non-reasoning model
families from three commercial API providers and one open-weight family:
OpenAI's GPT-4o (2024-08-06) \citep{openai2024gpt4o}, which serves as the
anchor model present in every cell; Anthropic's Claude 3.5 Haiku (2024-10-22)
\citep{anthropic2024claude}; Google's Gemini 2.0 Flash (2025-02-05)
\citep{google2024gemini}; and Google's open-weight Gemma 27B
\citep{google2024gemma}. Two further models delimit the capability range in
Appendix~\ref{app:ablations}: gpt-5-mini (2025-08-07), a reasoning model that
serves as the frontier tier, and Meta-Llama-3-8B-Instruct, a small
open-source model that serves as the capability floor. Our goal is to study the
interaction between cognitive constraints and mechanism design, which requires
models that are capable enough to parse the rules of a mechanism yet still
susceptible to the reasoning failures that mechanism design aims to remedy. We
therefore deliberately select well-established, widely deployed model families
rather than the most powerful frontier reasoning models: at the frontier, raw
reasoning ability may mask the very cognitive bottlenecks our levers target,
collapsing our ability to identify which design features matter and why. The
frontier is a measured scope condition, not an omission---and it is
two-sided. In the fidelity battery, gpt-5-mini plays the textbook formats
essentially perfectly (SPSB-IPV SMAD $0.2\%$; FPSB bids within pennies of the
risk-neutral equilibrium shading rule), yet it misses the third-price
format's equilibrium by a wide margin (Appendix~\ref{app:ablations}): frontier
reasoning narrows, but does not uniformly dissolve, the bounded-rationality
regime in which the ranking applies.

\paragraph{Temperature.}
All non-reasoning models are run at temperature $T=0.5$, the most common
setting in LLM-based simulation \citep{horton2023large,Horton2024EDSL}; for
reasoning models (gpt-5-mini), temperature is ignored by the API. For the
anchor model we ablate $T\in\{0.1, 1.0\}$: effects are modest and
non-monotonic across formats, so $T=0.5$ is retained as a comparability
convention rather than a performance choice
(Appendix~\ref{app:ablations}, Table~\ref{tab:temperature}).

\paragraph{One-shot protocol and the learning ablation.}
Our main experiments use \emph{single-instance} simulations: each repetition
$k$ is an independent market instance with freshly drawn values or preferences,
and agents do not receive feedback across instances. We adopt this design to
minimize reliance on cross-instance learning and long-term adaptation, which
remains unreliable for language models \citep{dong2024survey}. This choice is
distinct from whether a mechanism is \emph{iterative within} an instance: clock
auctions and iterative DA still unfold over multiple within-instance steps, but
we do not run repeated markets with feedback by default. In
Appendix~\ref{app:learning}, we run a repeated-auction ablation in which
bidders observe outcomes from prior instances; consistent with the motivation
above, we find no meaningful learning effects across repetitions, which also
justifies pooling one-shot observations across the two source studies.

\paragraph{Simulation protocol.}
Each experiment proceeds as follows, repeated for $K$ independent repetitions:
\begin{enumerate}
    \item \textbf{Market realization.} For repetition $k$, we draw values
    (IPV/APV), signals (CV), or school-choice valuations (DA) for all agents
    under a fixed random seed.
    \item \textbf{Prompted action.} We construct an individualized prompt for
    each agent containing (i) the mechanism rules (and treatment text, if
    applicable), (ii) the agent's private value, signal, or preference
    information, and (iii) in the learning ablation only, a feedback transcript
    summarizing prior instances. We request a structured response with explicit
    tags: \texttt{<PLAN>}\,\ldots\,\texttt{</PLAN>} /
    \texttt{<ACTION>}\,\ldots\,\texttt{</ACTION>} in auctions and
    \texttt{<REASON>}\,\ldots\,\texttt{</REASON>} /
    \texttt{<DECISION>}\,\ldots\,\texttt{</DECISION>} in DA, where the action
    tag contains a numeric bid (sealed bid), a stay/exit decision (clock), an
    increase/hold decision (eBay), a rank-order list (direct DA), or a yes/no
    answer (iterative DA).
    \item \textbf{Parallel elicitation and parsing.} We query the LLM for each
    agent using separate API calls in parallel, and parse the action from the
    action tag using a deterministic regex-based extractor.
    \item \textbf{Outcome computation.} Given the submitted actions, we
    determine allocations and payments mechanically---winner and price under
    the auction rules, or the final matching by executing student-proposing DA
    (or the query protocol) with the true school priorities.
\end{enumerate}
Unless otherwise noted, auction settings use $n=3$ bidders, and each
configuration is repeated $K$ times with independent seeds
(Table~\ref{tab:calibration-master}). Step-by-step procedures for the
sealed-bid, clock, and eBay environments appear in
Appendix~\ref{app:procedures}.

\subsection{Metrics}\label{sec:metrics}

The ranking requires a single deviation scale that is comparable across
mechanisms, environments, and models. We build it in three steps: a scaled
deviation metric for bids, its analogue for rank-order reports, and a
preservation index that expresses every lever's effect relative to its own
baseline.

\paragraph{Unified deviation metric for auctions (applied to all bids).}
Consider an auction \emph{environment} and \emph{format} indexed by $m$. The
index $m$ collects all primitives needed to define the theoretical benchmark,
including the number of bidders $N$, the value (or signal) distribution, the
information structure (independent private value/common value), and the
mechanism rules (e.g., first-price vs.\ second-price, sealed-bid vs.\ clock).
Let $\mathcal I_i$ denote bidder $i$'s information at the moment of bidding
under environment $m$ (e.g., $\mathcal I_i=v_i$ in private values, or
$\mathcal I_i=s_i$ in common values). Let $b_{i}$ denote an \emph{observed
submitted bid} by bidder $i$ (we apply the definition to \emph{all bids}, not
only winning bids), and let $b_m^\star(\mathcal I_i)$ denote the
\emph{theoretical benchmark bid} under $m$ given information $\mathcal I_i$:
the risk-neutral Bayes--Nash equilibrium in non-strategy-proof formats,
\[
b^\star_{\text{FPSB}}(v) = \tfrac{n-1}{n}\, v, \qquad
b^\star_{\text{TPSB}}(v) = \tfrac{n-1}{n-2}\, v,
\]
and truthful bidding, $b^\star(v)=v$, in strategy-proof formats (second price,
ascending clock).

\noindent We define the absolute bid error for a bid with information
$\mathcal I_i$ as
\[
d_m(\mathcal I_i)\;\equiv\;\bigl|b_i - b_m^\star(\mathcal I_i)\bigr|.
\]
To compare deviations across formats on a common scale, we normalize by the
expected benchmark bid level,
\[
\mu_m^\star \;\equiv\; \mathbb E_m\!\left[b_m^\star(\mathcal I)\right],
\]
where $\mathbb E_m[\cdot]$ denotes expectation with respect to the distribution
of the information object $\mathcal I$ induced by environment $m$ (i.e., the
value/signal distribution and any randomization in the experimental design). We
then define the \emph{scaled mean absolute deviation} (SMAD) for auction format
$m$ as
\[
\Delta_m \;\equiv\; 100 \cdot \frac{\mathbb E_m\!\left[d_m(\mathcal I)\right]}{\mu_m^\star}
\;=\; 100 \cdot \frac{\mathbb E_m\!\left[\left|b-b_m^\star(\mathcal I)\right|\right]}{\mathbb E_m\!\left[b_m^\star(\mathcal I)\right]}.
\]
This metric is interpretable as a percentage of the expected benchmark bid:
$\Delta_m=10$ means that, on average, bids deviate from the benchmark by an
amount equal to $10\%$ of the expected benchmark bid.

Given a sample of bids $\{(b_k,\mathcal I_k)\}_{k=1}^n$ from format $m$, the
plug-in estimator is
\[
\widehat{\Delta}_m
\;=\;
100 \cdot
\frac{\frac{1}{n}\sum_{k=1}^{n}\left|b_k-b_m^\star(\mathcal I_k)\right|}
{\widehat{\mu}_m^\star},
\qquad
\widehat{\mu}_m^\star \equiv \frac{1}{n}\sum_{k=1}^{n} b_m^\star(\mathcal I_k).
\]
When we report a standard error for the (unscaled) mean absolute deviation
$\frac{1}{n}\sum_k |b_k-b_m^\star(\mathcal I_k)|$, we propagate it to
$\widehat{\Delta}_m$ by scaling with $\widehat{\mu}_m^\star$. We construct
$95\%$ confidence intervals using a normal approximation when an analytic
standard error is available. Alongside SMAD, we report the \emph{signed} mean
deviation ($\text{bid}-\text{benchmark}$, in dollars) where the direction of
error is itself of interest; the two statistics can disagree about model
orderings, and we treat that disagreement as informative rather than choosing
between them.

\paragraph{Deterministic benchmarks.}
SMAD requires a single-valued benchmark mapping
$\mathcal I \mapsto b_m^\star(\mathcal I)$. Accordingly, we restrict attention
to auction formats where the theoretical prediction is deterministic given the
bidder's information.

\paragraph{Matching analogue: normalized Kendall-$\tau$ error.}
For matching we construct the same kind of object: a unit-free percentage
deviation from truthful play. For student $i$ with true preference ranking
$\pi_i^\star$ and submitted ranking $\pi_i$ over the school set $S$
($|S|=4$), the \emph{normalized Kendall-$\tau$ error} is the fraction of
school pairs ranked in opposite order,
\[
\kappa_i \;\equiv\; 100 \cdot
\frac{\#\bigl\{\{s,s'\} : \pi_i \text{ and } \pi_i^\star \text{ disagree on } s \text{ vs.\ } s'\bigr\}}
{\binom{|S|}{2}},
\]
so $\kappa_i = 0$ if and only if the report is truthful and $\kappa_i = 100$
for a complete reversal. The parallel to SMAD is deliberate but not literal:
SMAD normalizes the deviation by the expected benchmark bid, while $\kappa$
normalizes by the maximum number of pairwise inversions. Both are unit-free
percentages that equal zero exactly at truthful play and increase in the size
of the deviation, which is what the ranking requires. Because $17\%$ of
simulated students draw at least one pair of exactly tied school values, we
also compute a tie-robust variant of $\kappa$ that does not count inversions
among tied schools; headline DA numbers use the tie-robust variant, since
apparent ``misreports'' that merely permute tied schools are artifacts of the
deterministic tie-break (Appendix~\ref{app:da}). Even small $\kappa$ can
be consequential: under deferred acceptance, a single misreport can cascade
through rejection chains and alter the final matching for $\Theta(n)$
participants \citep{dubins1981machiavelli}.

\paragraph{Misreports under the iterative protocol.}
The iterative DA protocol never elicits a full rank-order list, so $\kappa$
cannot be computed directly. Instead we classify each revealed yes/no answer
against the student's true preferences: a \emph{Type-1 misreport} (false
rejection) answers NO when the queried school is in fact the most preferred
among the remaining set; a \emph{Type-2 misreport} (false acceptance) answers
YES when a more preferred school is still available. The iterative error rate
is the share of misreports among revealed decisions. Because applicants may
not be asked to reveal their entire rank-order lists, we count misreports from
the smaller set of revealed preferences: a measured error rate of zero
therefore bounds only the \emph{observable} component of behavior---models
could still be making cognitive errors that never surface because the protocol
never queries them. This observability caveat travels with every use of the
iterative-DA result in this paper. Formal per-node definitions, the per-model
decision-level accounting, rule-of-three upper confidence bounds for cells
with zero observed misreports, and a comparison of the pick-based and yes/no
query implementations appear in Appendix~\ref{app:da}.

\paragraph{Baseline declaration.}
Two baseline families appear in the auction analysis, and we fix their use
once. \emph{Cross-format} comparisons---sealed bid versus clock, the OSP rung
of the ranking---are made against the dedicated SPSB baseline cells, whose
mean deviations ($-0.49$, $-1.63$, $-3.28$, $-5.27$ for Claude 3.5 Haiku,
Gemini 2.0 Flash, GPT-4o, and Gemma 27B) match the predecessor draft's
published values to rounding. \emph{Within-format} description and scaffold
contrasts (Families B--D) are evaluated against the same model's axis-specific
baseline cells run in the same batch (pooled across axes where a treatment
spans them), which differ materially from the dedicated cells for some models
(e.g., $+0.48$ versus $-0.49$ for Claude 3.5 Haiku). Both baseline families
are reported in the released cell-level data, and no comparison in this paper
mixes them.

\paragraph{Preservation index.}
For lever $L$, domain $d\in\{\text{auction},\text{DA}\}$, and model $m$, let
$D(L,d,m)$ denote the mean scaled deviation from dominant-strategy play---SMAD
for auctions, mean normalized Kendall-$\tau$ error for DA---and let
$D(\text{base},d,m)$ denote the same object under the domain's baseline
mechanism and description (SPSB sealed bid; direct DA). The
\emph{preservation index} is
\[
\rho(L,d,m) \;\equiv\; 1 - \frac{D(L,d,m)}{D(\text{base},d,m)}.
\]
By construction $\rho=1$ means the lever restores perfect dominant-strategy
play, $\rho=0$ means it has no effect, and $\rho<0$ means it \emph{backfires},
pushing behavior further from the dominant strategy than the untreated
baseline. Because $\rho$ is computed within each model$\times$domain cell, the
ranking it induces is invariant to cross-model differences in baseline levels
(and to the choice between absolute and signed deviation measures for ordering
models)---a lever is credited only for what it does relative to the behavior
it inherits.

\paragraph{Tier-claim discipline.}
We pre-commit to reporting the ranking as a \emph{tier structure} (a partial
order over levers), not a strict total order. Adjacent tiers are separated
only where pairwise sign tests across the model$\times$domain cells support
the separation; concordance of the full ordering across cells is summarized by
Kendall's $W$ with bootstrap confidence intervals; and within-tier reversals
across domains (e.g., contingent scaffolds versus safety descriptions) are
reported as findings, not smoothed away. Section~\ref{sec:ranking-synthesis}
implements this discipline.

\paragraph{Statistical standard.}
Every reported cell receives a bootstrap standard error and 95\% confidence
interval (resampling market instances, $2{,}000$ replications, seed 1299), and
every treatment effect is evaluated against the baseline declared above with
both a Mann--Whitney test and a Welch $t$-test, along with Cohen's $d$; we
report both tests because description levers can move dispersion without
moving the mean (and vice versa), and the two tests can disagree in exactly
those cases. Where an error rate of exactly zero is observed, we report
one-sided rule-of-three upper confidence bounds rather than declaring
perfection (Appendix~\ref{app:da}). Concordance of the lever ordering across
model$\times$domain cells (Kendall's $W$ with bootstrap confidence intervals),
adjacent-tier sign tests, and the double-bootstrap intervals for the pooled
preservation indices are reported in Section~\ref{sec:ranking-synthesis}.
Knife-edge cells (e.g., near-zero error rates under Rejection Safety and
iterative DA) are topped up to $K=100$ before adjacent rungs are declared
separated.

\subsection{Human Benchmark Reconstruction}\label{sec:reconstruction}

\paragraph{Reconstructing human bid distributions via moment matching.}
Because micro-level bid data are rarely available from published experiments,
we reconstruct synthetic human bid distributions by calibrating a mixture model
to aggregate statistics reported in each source. We model bids as:
\[
b = b^\star_m(v) \cdot r
\]
where $b^\star_m(v)$ is the equilibrium bid and $r$ is a bid-to-equilibrium
ratio drawn from a three-component mixture:
\begin{align*}
\text{With probability } p_{\text{eq}}: \quad & r \sim \mathcal{N}(1, \sigma^2) \\
\text{With probability } p_{\text{over}}: \quad & r \sim 1 + \text{Exp}(\lambda) \\
\text{With probability } p_{\text{under}}: \quad & r \sim \text{Uniform}(1-\delta, 1)
\end{align*}
This structure maps directly to how experimental papers categorize behavior:
equilibrium bidders ($r \approx 1$), overbidders ($r > 1$), and underbidders
($r < 1$).

\paragraph{Source papers and equilibrium benchmarks.}
We draw on four sources: \citet{li2017obviously} and
\citet{breitmoser2022obviousness} for second-price and ascending clock
auctions; \citet{kagel1993independent} for first-price, second-price, and
third-price sealed-bid auctions with $n \in \{5, 10\}$ bidders; and
\citet{gonczarowski2023strategyproofness} for second-price auctions comparing
mechanism descriptions. For each format, we use the risk-neutral Bayes--Nash
equilibrium benchmarks of Section~\ref{sec:metrics}.

\paragraph{Parameter identification.}
Mixture weights $(p_{\text{eq}}, p_{\text{over}}, p_{\text{under}})$ are set
directly from reported rates. For example, \citet{li2017obviously} reports
${\sim}50\%$ dominant-strategy play and ${\sim}40\%$ overbidding, yielding
$p_{\text{eq}} = 0.50$, $p_{\text{over}} = 0.40$, $p_{\text{under}} = 0.10$.
The overbid scale $\lambda$ is calibrated to match a secondary moment:
\begin{itemize}
  \item For \citet{li2017obviously} and \citet{breitmoser2022obviousness}: the
  mean bid-to-value ratio (e.g., $\mathbb{E}[b/v] \approx 1.08$).
  \item For \citet{kagel1993independent}: the regression coefficient $\beta$
  from bid regressions.
  \item For \citet{gonczarowski2023strategyproofness}: mean absolute deviation
  (MAD $\approx$ \$0.51--0.55).
\end{itemize}
The equilibrium noise $\sigma$ is derived from regression $R^2$ where reported,
via $\sigma^2 = \beta^2 \text{Var}(v)(1-R^2)/R^2$. The underbid bound $\delta$
is set consistently across papers; this is a parsimony choice as papers do not
report underbid tail shapes separately. Full parameter values appear in
Appendix~\ref{app:mixture} (Table~\ref{tab:mixture-calibration}).

\paragraph{Why a mixture model rather than QRE.}
Quantal response equilibrium predicts symmetric deviations around equilibrium,
contradicting the well-documented asymmetric overbidding in second-price
auctions \citep{kagel1993independent}. The mixture model directly maps to how
experimental papers categorize and report behavior.

\paragraph{Limitations and epistemic status.}
Some reported statistics are difficult to reconcile within any parametric
model---for instance, the regression coefficients and bidding frequencies in
\citet{kagel1993independent} cannot both be matched under Gaussian noise,
suggesting heteroskedasticity in the original data. Our approach matches key
aggregate patterns while acknowledging this limitation. More fundamentally, the
reconstructions are model-based simulations of published moments, not samples
of human behavior, and we hold every use of them in this paper to a
corresponding standard: \emph{comparisons to reconstructed humans are made at
the level of rankings, directions, and comparative statics only; levels and
significance tests are reserved for theory benchmarks.} We run no hypothesis
tests against reconstructed distributions---such tests would be statistically
circular, reflecting our calibration choices rather than sampling variation in
human behavior---and where reconstructed anchors appear in figures, their
uncertainty is represented by Monte-Carlo bands obtained by re-drawing the
mixture parameters within the ranges the source statistics admit. The ranking
in Section~\ref{sec:ranking} never depends on the reconstructions; they enter
only in Section~\ref{sec:validation} (fidelity) and Section~\ref{sec:humans}
(anchoring). Full reconstruction details, per-paper calibrations, and the
reconstruction-uncertainty procedure appear in Appendix~\ref{app:human}.
